% 量子力学1 量子力学1-11 問題1（081-box） 解いた日 2026-11-20
%% kotae
\toi{1}{◎}{}
\[ E_{n}=\frac{n^{2}\pi^{2}\hbar^{2}}{2mL^{2}} \]
\[ E_{2}-E_{1}=\frac{3\pi^{2}\hbar^{2}}{2mL^{2}}=\frac{3h^{2}}{8mL^{2}} \]
\[ E_{2}-E_{1}=\frac{hc}{\lambda_{21}}\qquad\therefore\ \frac{3h^{2}}{8mL^{2}}=\frac{hc}{\lambda_{21}} \]
\begin{align*}
L^{2}&=\frac{3h\lambda_{21}}{8mc}\\
L&=\sqrt{\frac{3h\lambda_{21}}{8mc}}\\
&=\qty{6.61e-10}{\metre}\\
&=\qty{0.661}{\nano\metre}
\end{align*}
\toi{2}{△}{未着手}
\begin{align*}
E_{1}&=\frac{\pi^{2}\hbar^{2}}{2mL^{2}}\\
&=\qty{0.861}{\electronvolt}
\end{align*}
\nuke{$\lambda_{32}=hc/(5E_{1})=\qty{288}{\nano\metre}$}
\naze{問われているもう1つの量 $\lambda_{32}$ が書かれていない。$E_{n}\propto n^{2}$ なので $E_{3}-E_{2}=(9-4)E_{1}=5E_{1}$ で，波長は準位の差に反比例するから $\lambda_{32}=hc/(5E_{1})=(3/5)\lambda_{21}=\qty{288}{\nano\metre}$ である。}
\toi{3}{◎}{}
i) $n=1$
\begin{align*}
\int_{0}^{\frac{L}{4}}\abs{\phi_{1}(x)}^{2}\dd x&=\int_{0}^{\frac{L}{4}}\frac{2}{L}\sin^{2}\frac{\pi x}{L}\\
&=\frac{2}{L}\int_{0}^{\frac{L}{4}}\frac{1}{2}\left(1-\cos\frac{2\pi x}{L}\right)\dd x\\
&=\frac{1}{L}\left[x-\frac{L}{2\pi}\sin\frac{2\pi x}{L}\right]_{0}^{\frac{L}{4}}\\
&=\frac{1}{L}\left(\frac{L}{4}-\frac{L}{2\pi}\sin\frac{\pi}{2}\right)\\
&=\frac{1}{4}-\frac{1}{2\pi}
\end{align*}
ii) $n=2$
\begin{align*}
\int_{0}^{\frac{L}{4}}\abs{\phi_{2}(x)}^{2}\dd x&=\int_{0}^{\frac{L}{4}}\frac{2}{L}\sin^{2}\frac{2\pi x}{L}\dd x\\
&=\frac{2}{L}\int_{0}^{\frac{L}{4}}\frac{1}{2}\left(1-\cos\frac{4\pi x}{L}\right)\dd x\\
&=\frac{1}{L}\left[x-\frac{L}{4\pi}\sin\frac{4\pi x}{L}\right]_{0}^{\frac{L}{4}}\\
&=\frac{1}{L}\left(\frac{L}{4}-\frac{L}{4\pi}\sin\pi\right)\\
&=\frac{1}{4}
\end{align*}
$n=1$ のときは $1/4$ よりも小さい確率になる。

$n=2$ のときは古典的な粒子と一致する。
%% kizuki
\item \mae{第7章 問題3・11-04，第9章 問題7・11-11} 問いの一部を答えないまま終えた。(2) は $E_{1}$ と $\lambda_{32}$ の2つを求める問いである。解き終えたら問題文の「求めよ」の数を数え，答えた数と合わせる。
\item 無限井戸の準位の差はどれも $E_{1}$ の整数倍（$E_{n}-E_{m}=(n^{2}-m^{2})E_{1}$）なので，遷移の波長どうしの比は差の逆比で，$L$ を使わずに求まる。上の準位ほど間隔が広いので $\lambda_{32}<\lambda_{21}$（$5/3$ 倍ではなく $3/5$ 倍）で確かめる。
